\section{Theorem V3.1 --- donor conservativity}

Adding the scientific-discharge layer does not change any donor-native action, delegation, evidence, adjudication, or release verdict. Local authorities remain usable for their native purpose even when they do not discharge a broader scientific target.

\section{Theorem V3.2 --- cross-donor scientific-discharge separation}

For every registered donor family and each non-inert coordinate of \texttt{(domain, kind, scope, content, epoch)}, native donor success can coexist with scientific-discharge failure when only that coordinate is incompatible with the target obligation.

Thus a valid action certificate, verified claim-evidence packet, scientific release decision, or cross-domain authorization may be retained without being treated as an unrestricted authority token.

\section{Theorem V3.3 --- monotone scientific authority propagation}

Scientific authority may be preserved or narrowed through heterogeneous authority chains. A hop that widens the scientific type beyond inherited authority cannot discharge the widened target merely because each local donor step remains native-valid.

This extends donor narrowing/cascade-containment ideas to the scientific-obligation type rather than replacing them.

\section{Theorem V3.4 --- protected coercion}

A complete protected coercion may transform an otherwise incompatible scientific type. The same mismatch without the registered bridge remains blocked. Semantic similarity, intent similarity, common action digest, claim fluency, or ordinary action composability is not sufficient to create a coercion.

\section{Theorem V3.5 --- three-state blockers}

Under otherwise satisfied conditions:

\begin{itemize}
  \item \texttt{REFUTED} permits the blocker condition to clear and allows discharge;
  \item \texttt{UNDETERMINED} yields \texttt{CANNOT_CHECK};
  \item \texttt{ESTABLISHED} blocks discharge.
\end{itemize}

The distinction prevents unknown blocker state from being silently treated as blocker absence.

\section{Theorem V3.6 --- exact support-family revocation}

When two independent complete support families establish the same target obligation, revoking one family preserves discharge through the other. Revoking all complete support families removes discharge.

This avoids both over-revocation, in which one failed premise destroys an independently supported result, and under-revocation, in which a broken unique support path is ignored.

\section{Theorem V3.7 --- heterogeneous authority-chain composition}

Every ordered pair among the thirteen registered donor families admits a scientifically valid two-hop composition under native-valid local verdicts, compatible or protectedly bridged scientific type, non-widening authority, refuted blockers, and surviving support. The matched chain with a scientific-authority-widening hop is blocked even though both donor-native steps remain valid.

The theorem treats heterogeneous authority composition as a first-class object rather than requiring every donor to adopt one common receipt format.

\section{Theorem V3.8 --- permission, release, and support are not inverse relations}

A native denial of one action or release route is not scientific refutation when an independent complete evidence family establishes the proposition. Conversely, native authorization or release does not authorize a different or broader scientific target. The two relations are connected by explicit evidence/type bridges rather than logical inversion.

\section{Theorem V3.9 --- heterogeneous receipt composition is a prerequisite, not scientific sufficiency}

Valid composition of authorization receipts, canonical action or subject binding, and successful local policy evaluation may be necessary inputs to a scientific claim, but they do not by themselves establish target scientific sufficiency. Scientific discharge remains governed by the target type, blocker, support, and coercion relation.

\section{Theorem V3.10 --- ideal decentralized-product equivalence}

A decentralized donor product supplied with the exact same scientific type, narrowing, blocker, support-family, coercion, and composition rules agrees extensionally with P8. The contribution is therefore the explicit common scientific-discharge calculus, not centralized expressive power.
